\documentclass[11pt]{article}
\usepackage[margin=1in]{geometry}
\usepackage{amsmath,amssymb,amsthm,booktabs,microtype,hyperref}
\hypersetup{hidelinks,pdfauthor={T},pdftitle={Matched Preference Labels Under Source Failures}}
\newtheorem{proposition}{Proposition}
\title{Matched Preference Labels Under Source Failures}
\author{T}
\date{}
\begin{document}
\maketitle

\begin{abstract}
Matched preference labels should agree or reverse under specified
changes to a scenario. Signed-graph balance detects inconsistent
rules and gives a cycle that an editor can inspect. When reference
labels may share a faulty source, a vertex-cover calculation
identifies which label orientations are still forced. Exhaustive
finite checks verify both reductions. A separate lexical experiment
on public deontology annotations compares responses within the same
scenario, tests the hard split, and measures how often the same
responses change ranking across scenarios. The structural guarantee
requires valid matching rules and a source-fault budget.
\end{abstract}

\section*{Summary and scope}
\paragraph{Question.} What can paired supervision establish when
comparisons are consistent but the reference labels are uncertain?
\paragraph{Finding.} A signed graph certifies consistency. A separate
vertex-cover test certifies which orientations remain forced under
a stated source-fault budget.
\paragraph{Evidence and scope.} The reductions are proved
and checked with exact oracles. A fitted lexical study on a public
corpus evaluates matching as a diagnostic, without attesting the
correctness of an annotation, a person's preferences, or an agent's
future actions.

\section{What a swap should preserve}
Counterfactual fairness treats certain demographic changes as
interventions whose effect on a decision depends on a causal
model \cite{kusner}. During data construction the narrower task is
to enforce a declared symmetry only when relevant roles, permissions
and facts are held fixed. Swapping a name under unchanged authority
may leave the preferred action the same. Swapping the roles of
the person who gave consent and the person who did not can legitimately
change it. A pair needs a reason for its sign.

Label a finite set of matched cases by $y_i\in\{-1,1\}$,
meaning one of two clearly ordered responses wins. Each edge
requests $y_i y_j=s_{ij}$, where sign $+1$ preserves the preference
and sign $-1$ reverses it. The graph-balance characterization is the classical criterion of
Harary, as presented by Zaslavsky \cite{zaslavsky}. For example,
if two successive swaps preserve a label, the direct comparison
between their endpoints cannot reverse it. The three edges form a
negative cycle. An editor can inspect those three transformation
rules before deciding whether to collect more labels.

\begin{proposition}
The constraints are satisfiable exactly when every cycle has
positive sign product.
\end{proposition}
\begin{proof}
Multiplying the constraints around a cycle cancels each $y_i$
twice, giving one. Conversely fix a root label in each component
and propagate labels by multiplying along a spanning tree.
Every non-tree edge closes a cycle. Its sign equals the product
along the tree path by positivity of that cycle, so its constraint
is satisfied.
\end{proof}

\section{What the developer should do with a conflict}
When a cycle fails, preserve the source cases and the transformation
that generated each edge. Review the first inconsistent cycle before
tuning the model. A missing role distinction can justify deleting
or reversing an edge with a logged reason. Ambiguous consent may
need a third response type, which lies outside the binary model.
Anchor valid components with independently adjudicated examples.
Without an anchor, flipping every label in a component preserves
all edges, so structural consistency cannot identify which action
is actually appropriate.

\section{A shared fault budget for anchor sources}
After the balance check, fix one satisfying labeling $g_i$ on
each component. Every other satisfying labeling has the form
$y_i=x_{C(i)}g_i$, where $x_c\in\{-1,1\}$ is the orientation of
component $c$. An anchor from root $r$ asserting $y_i=b$
endorses $x_{C(i)}=bg_i$. All reports traced to the same original
adjudication are charged to the same root, even if they arrive
through different reviewers. Assume the signed relations are
correct and that some set of at most $f$ roots accounts for
every incorrect anchor. No probabilistic independence assumption
is needed.

Let $F$ contain the roots that endorse both signs of any one
component. Form the conflict graph $H$ on the other roots,
joining two roots when their endorsements disagree on a
component. For a partial orientation $z$, let $F_z$ be $F$
together with every root that contradicts any of its pinned
orientations. Write $\tau(G)$ for the minimum size of a vertex
cover of $G$.

\begin{proposition}
A partial orientation $z$ extends to a labeling compatible with
the fault budget exactly when
\[
 |F_z|+\tau(H-F_z)\le f .
\]
\end{proposition}
\begin{proof}
Every root in $F_z$ must be removed. Every edge left in $H-F_z$
has incompatible endpoints, so the remaining removals must cover
all its edges. Conversely, remove $F_z$ and a cover of those
edges. Each component now receives either one sign or no sign
from the surviving roots. Its sign can be set to that unanimous
endorsement, to its pin if it has one, or freely in the absence
of both. The resulting labeling satisfies all the signed pairs
and all surviving anchors.
\end{proof}

Test both pins for each component. If only one is feasible, the
budget forces it. If both are feasible, retain the ambiguity.
If neither is feasible, the budget and reports are jointly
inconsistent. For the bounded cover search, choose a remaining
edge and branch on removing either endpoint. A residual budget
of $k$ therefore requires at most $2^{k+1}-1$ search nodes,
with polynomial work at each node. For unbounded $f$ the
feasibility problem is NP-complete by the classical vertex cover
problem \cite{karp}. To realize an arbitrary conflict graph,
introduce one singleton component per edge and have the two
endpoint roots endorse opposite signs on it. A set of removed
roots is a certificate whose consistency is checked directly.

Suppose root A endorses the positive sign on two components,
while roots B and C each endorse the negative sign on a
different one. A budget of one leaves either sign possible
when each component is examined alone. Jointly, the only
one-root cover is $\{\mathrm{A}\}$, and both orientations are
negative. If A's reports are copied through more reviewers,
the conclusion is unchanged when provenance is preserved.
An editor can use this test to direct adjudication toward
components that remain ambiguous. If the signed comparisons,
source identities, or fault budget are wrong, the result
does not certify the resulting labels.

\section{Exhaustive check and model experiment}
Each pair of four vertices can be absent, positive, or negative.
The audit compares traversal with all sixteen complete assignments
for $3^6=729$ signed graphs. Both methods classify 393 as
consistent. The three-edge witness with two positive and one
negative edge has no solution. For the source-fault extension,
direct orientation enumeration agreed with the cover solver on
472,392 pinned queries across 19,683 ternary matrices and four
budgets. In 9,792 component-and-budget cases a feasible global
model forced a sign that the isolated component left ambiguous.
The solver also agreed with exhaustive source-subset search on
all 1,024 five-root conflict graphs, with six budgets per graph.
These counts describe the finite test design. A later study should
inject matched-example contradictions and compare training after
source-level review with training after discarding examples
without explanation. Evaluate justified asymmetries separately from
identity invariance. Consistency of declared pairs still depends on
their causal and normative validity.
\section{Learning from independently anchored pairs}
We generated paired supervision with an unchanged task signal and
a swapped identity cue. In a designated quadrant we corrupted
90\% of positive labels. The repaired training arm used the
simulator's independent truth for each pair as its anchor and
propagated the known preserving sign to the other member.
Both arms fitted logistic scorers of the same form. On 3,200
held-out pairs, the raw labels produced 1,646 swap disagreements
and 4,740 correct individual decisions out of 6,400. Anchored
labels yielded zero disagreements and 6,400 correct decisions.
This experiment models valid anchor information. Agreement
alone would also allow both members of a pair to be wrong.

\section{Matched scoring on public deontology judgments}
The ETHICS deontology files pair scenarios with responses
and include binary judgments \cite{hendrycks}. They provide
multiple responses to the same scenario. This experiment
asks how well a lexical scorer orders those responses.
It does not infer signed invariances or adjudicate the
provenance of individual labels.

Before fitting, we fixed four feature views and hashed the
data. We fit add-one lexical likelihood ratios from 18,161
training rows after removing three duplicates. Scenario
words, response words and a small relation block have
separate likelihood ratios. Each vocabulary keeps features
seen at least three times. The relation block includes
shared tokens, a capped overlap count and joint negation
flags. Compare scenario alone, response alone, their
additive sum, and the additive sum with the relation block.
The same training class prior is used for every row
classification decision. There is no test tuning.

We exclude test scenarios whose case-and-whitespace-normalized
text appears in training. This removes 1,300 ordinary rows and
736 hard rows. The remaining splits share one normalized
scenario with each other. This filter addresses exact scenario
overlap and does not test for paraphrases. Within a scenario,
compare every label-one response with every label-zero response,
award half credit for ties, and average the resulting proportions
equally across scenarios. There are 2,252 comparisons across
574 ordinary scenarios and 2,742 across 700 hard scenarios.

\begin{center}
\small
\begin{tabular}{lrrrr}
\toprule
 & \multicolumn{2}{c}{Ordinary, 2,296 rows}
 & \multicolumn{2}{c}{Hard, 2,800 rows} \\
Score & Row correct & Pair order & Row correct & Pair order \\
\midrule
Scenario & 1,160 & 0.5000 & 1,402 & 0.5000 \\
Response & 1,618 & 0.8077 & 1,696 & 0.6597 \\
Additive & 1,580 & 0.8077 & 1,693 & 0.6597 \\
Additive with relations & 1,569 & 0.8046 & 1,693 & 0.6596 \\
\bottomrule
\end{tabular}
\end{center}

The scenario-only score ties by construction. In the additive
score the scenario term cancels between two responses to the
same scenario, so its ranking is exactly the response-only
ranking. Its row decisions can still change. The tested
relation features failed to improve within-scenario ranking
on either split. The response-only score lost about 14.8
percentage points on the hard split. These two splits do
not identify which cues caused that change. A dataset editor
can use this matched diagnostic alongside row scores to check
whether a model distinguishes plausible responses to the
same request, while auditing label validity separately.

\section{When the scenario must change the answer}
Matching controls the scenario within each comparison. A benchmark
can still reward a response-only score if the same responses seldom
need different treatment across scenarios. For normalized response
$e$, let $n_{e,0}$ and $n_{e,1}$ count its observed labels. Every
deterministic response-only classifier makes at least
$\sum_e\min(n_{e,0},n_{e,1})$ row errors. Predicting the empirical
majority for each response achieves that minimum on the given rows.
That oracle uses the test labels and serves as an identifiability
check, without estimating a deployable classifier's accuracy.

For rankings, let $w_{ab}$ be the total macro weight of comparisons
requiring response $b$ to score strictly above $a$. Each mixed
scenario has equal weight, split equally over its opposite-label
pairs. With half error for ties, any response-only score incurs
at least
\[
 \sum_{\{a,b\}}\min(w_{ab},w_{ba})
\]
macro error, where each unordered pair appears once. Ordering
the responses either way incurs one directional weight. Tying
them incurs half their sum. Each case costs at least the smaller
weight. An additive scenario term cancels within each pair, so
the same bound applies to additive scoring. Longer cycles can
force error even when every term in the sum vanishes. Exact
enumeration of weak orders verified the bound on 729 weighted
three-response configurations, with equality in 675.

In the ordinary split, no normalized response carries both labels.
In the hard split, six responses do, covering twelve rows and
forcing at least six errors out of 2,800 for a response-only rule.
Neither split contains an exact reversal of an observed response
pair. The direct-reversal floor is therefore zero on both. This
test exposes a small number of incompatible row labels. It cannot
tell whether each conflict reflects context, a missing detail, or
annotation error. It also cannot identify what caused the earlier
ranking gap.

There are 67 ordinary and 66 hard rows whose normalized response
appears in training. As an overlap sensitivity check, we dropped
every test scenario containing any such response, leaving 511
ordinary and 642 hard scenarios. With the scorer fixed, response
concordance remained $0.8037$ and $0.6635$, respectively. The
gap persisted under the exact-overlap filter. Paraphrase and template overlap remain
untested. A prospective benchmark could use adjudicated pairs of
scenarios in which the same two candidate responses trade places
under a controlled change in permission or relevant facts. The
resulting reversals would force a known error floor for any model
that ignores the scenario.

\paragraph{Review status.} Independent manuscript review and a language-model study remain open.

\begin{thebibliography}{9}
\bibitem{kusner} M. J. Kusner, J. Loftus, C. Russell and R. Silva.
Counterfactual Fairness. \url{https://arxiv.org/abs/1703.06856}.
\bibitem{zaslavsky} Thomas Zaslavsky. Signed graphs.
Discrete Applied Mathematics 4, 47--74, 1982.
\url{https://davinci.fmph.uniba.sk/~matok1/diplomovka/sources/Zaslavsky-former.pdf}.
\bibitem{karp} Richard M. Karp. Reducibility Among Combinatorial
Problems. In \emph{Complexity of Computer Computations},
85--103, 1972.
\url{https://cgi.di.uoa.gr/~sgk/teaching/grad/handouts/karp.pdf}.
\bibitem{hendrycks} Dan Hendrycks et al.
Aligning AI With Shared Human Values. \emph{ICLR}, 2021.
\url{https://arxiv.org/abs/2008.02275}.
\end{thebibliography}
\end{document}
